% ==================== 摘要 ====================
{\bfseries 摘要}\quad
红外与拉曼都探测分子振动，但物理来源与选择定则不同。本文以油酸包覆 TiO$_2$ 纳米颗粒为例，围绕“油酸以何种羧酸根配位、其与晶相有何关系”比较两种技术。本人认为应以红外判配位，拉曼补晶相与分布，联用可建立“配位—晶相—位置”的关联，并讨论水、荧光、制样与激光损伤等限制。

\vspace{4pt}

\par\noindent{\heiti 关键词：}红外光谱；拉曼光谱；TiO$_2$ 纳米颗粒；羧酸根配位；选择定则；联用
